\documentclass{article}
\usepackage[utf8]{inputenc}
\usepackage[english]{babel}
\usepackage[font=small,labelfont=bf]{caption}
\usepackage{geometry}
\usepackage[sort&compress, numbers, super]{natbib}
\usepackage{amsmath}
\usepackage{pxfonts}
\usepackage{graphicx}
\usepackage{newfloat}
\usepackage{setspace}
\usepackage{hyperref}
\usepackage{placeins}

\newcommand{\argmax}{\mathop{\mathrm{argmax}}\limits}

% Notation macros (kept in sync with main.tex)
\newcommand{\tauinit}{\tau}

\title{\textit{Supplementary materials for}: A Multi-Stream Temporal Context Model for narrative memory}
\author{Xinming Xu and Jeremy R.\ Manning\textsuperscript{*}\\
  Department of Psychological and Brain Sciences, Dartmouth College\\
  \textsuperscript{*}Corresponding author: \href{mailto:jeremy.r.manning@dartmouth.edu}{jeremy.r.manning@dartmouth.edu}
}

\begin{document}

\renewcommand{\figurename}{Supplementary Figure}
\renewcommand{\tablename}{Supplementary Table}

\setcounter{equation}{0}
\setcounter{figure}{0}
\setcounter{table}{0}
\setcounter{page}{1}
\setcounter{section}{0}
\makeatletter
\renewcommand{\theequation}{S\arabic{equation}}
\renewcommand{\thefigure}{S\arabic{figure}}
\renewcommand{\thetable}{S\arabic{table}}
\renewcommand{\bibnumfmt}[1]{[S#1]}
\renewcommand{\citenumfont}[1]{S#1}

\maketitle

\section*{S0. Architectural lineage of MS-TCM}

The present paper describes MS-TCM as a multi-stream extension of
\citet{CornellZhang2025} hierarchical CMR. The model evolved through
several preceding architectural revisions whose details are recorded
in the iteration log at
\texttt{notes/mstcm\_architecture\_log.md} and are not part of the
present specification. The earliest formulation
(\texttt{notes/ms-tcm.pdf}) used a single global context plus a bank
of storyline-specific context vectors mixed into an encoding composite
with fixed weights, plus a task-dependent retrieval reweighting; that
architecture failed to reproduce the canonical free-recall phenomena
(serial position curve, probability of first recall, lag-CRP) because
it bypassed the standard CMR retrieval route, and was retired in
favour of the v6 design (\texttt{notes/two\_level\_cmr\_v6.pdf})
which rebases MS-TCM on top of \citet{CornellZhang2025}. Subsequent
empirical observations in the FRFR-category worked example motivated
two further additions to the v6 design: a $\tauinit$ recall-onset
mixture (because the C\&Z 2025 baseline structurally cannot produce a
primacy spike in the probability of first recall) and a strict
hierarchical recall route with a separate global cross-storyline
context (because participants organize recall by category even when
the presentation order is randomized). The present paper consolidates
those additions into a single MS-TCM specification --- the model
described in \S2 of the main text and implemented at
\texttt{code/ms\_tcm/\_likelihood\_core\_mstcm.py}. The
iteration-by-iteration log of empirical motivations and architectural
decisions is at \texttt{notes/mstcm\_architecture\_log.md}; the symbol inventory and
file-level deletion manifest from the earlier-spec migration is at
\texttt{notes/v6\_migration.md}; git history preserves earlier
sources.

\section*{S1. Off-by-one correction in the reformatted recall table}

The upstream \texttt{exp2.egg} file from \cite{ManningEtAl2023FRFR} encodes
each recall event's provenance in a \texttt{temporal} HDF5 attribute that is
0-indexed against the \texttt{i\{k\}} subgroup names in the presented
group. The initial reformat (\texttt{scripts/reformat\_frfr\_category.py},
commit \texttt{8684fa0}) treated \texttt{temporal} as if it were 1-indexed,
which shifted every in-list recall's serial position down by one and
silently re-coded the last word of every list as an extra-list intrusion.
The corrected mapping (commit \texttt{405e272}, \texttt{sp = temporal + 1}
for \texttt{temporal} $\in [0, 15]$, else \texttt{sp = 0}) restores the
expected serial-position distribution (primacy at SP1--2, recency at
SP14--16) and reduces the extra-list intrusion count from $399$ to $10$.

A regression test at
\texttt{code/tests/test\_frfr.py::test\_recall\_serial\_position\_shape}
asserts $P(\text{recall}) > 0.3$ for the last three serial positions and
an intrusion rate below $5\%$, so any future off-by-one would fail CI.

\section*{S2. Likelihood-surface flatness on small synthetic data}

% TODO: re-run the grid-scan diagnostic against the current MS-TCM
% (beta_enc, gamma_fc) parameters and update the nat ranges accordingly.
% The original numbers below were computed for an earlier (beta_G, w_G)
% parameterization that no longer corresponds to the model and are kept
% here only as a placeholder for the qualitative pattern.
At FRFR-scale real data the MS-TCM likelihood surface is informative
(grid scans across the active parameters $(\beta_{\mathrm{enc}}, \gamma_{\mathrm{fc}})$ give a
range of $\sim 100$ nats across the parameter simplex, vs a
uniform-recall baseline $\sim 1200$ nats below the MLE). At the smaller
scales used for quick sanity checks, the same grid scan shows an
essentially flat surface (LL varies by $\sim 2$ nats around a uniform
baseline). The practical consequence is that parameter-recovery tests
need a synthetic dataset at FRFR scale ($30$ pt $\times$ $16$ lists
$\times$ $16$ words $\times$ $4$ categories $\approx 7{,}680$ recalls)
before the L-BFGS-B optimizer can reliably recover the true parameter
vector; smaller datasets converge to whatever corner the initial point
leans toward. This is a property of the likelihood surface, not a bug
in the fitter.

\section*{S3. Implementation details: norm-preserving drift, softmax gain, primacy, no-repeats}

Four implementation choices beyond the minimal equations in the main
text are necessary for the model's forward simulator and its
posterior-predictive to match the observed behavioural regularities
(SPC, P(first recall), lag-CRP, clustering). Each is standard in the
CMR / hierarchical-CMR literature but is worth stating explicitly.

\paragraph{Unit-norm feature encoder.}
Equation~(1) of the main text assumes unit-norm inputs so the drift
invariant $\|\mathbf{c}(t)\| = 1$ is preserved at every step. Raw
multi-hot feature vectors (with $k$ active slots) have norm
$\sqrt{k}$; the feature encoder therefore $L_{2}$-normalises each row
by default. See \texttt{code/ms\_tcm/features.py}
(\texttt{encode\_features(..., normalize=True)}).

\paragraph{Norm-preserving $\rho$ under non-orthogonal inputs.}
For strictly orthonormal inputs, \citet{HowaKaha02a}'s
$\rho = \sqrt{1 - \beta^2}$ preserves $\|\mathbf{c}(t)\| = 1$
exactly. Multi-hot feature vectors that share active slots across
items are unit-norm but not orthogonal, so the orthonormal formula no
longer maintains the invariant. We adopt the generalised formula
$\rho = \sqrt{1 + \beta^{2}(\dot{c}^{2} - 1)} - \beta\,\dot{c}$
(where $\dot{c} = \mathbf{c}_{\mathrm{prev}} \cdot \mathbf{c}^{IN}$)
from \citet{PolyEtal09} Eq.~3, which reduces exactly to
$\sqrt{1 - \beta^{2}}$ when $\dot{c} = 0$ and preserves unit norm for
arbitrary unit-norm inputs. This is applied at every drifting context
in MS-TCM: the per-storyline level (rate $\beta_{\mathrm{story}}$ on
storyline-$s$ items only), the global level (rate
$\beta_{\mathrm{enc}}^{G}$ on every step), and the retrieval-cue
drifts (rate $\beta_{\mathrm{rec}}$). See
\texttt{code/ms\_tcm/drift.py::\_norm\_preserving\_rho}.

\paragraph{Softmax gain $k$.}
The retrieval readout $a_{j} = M^{\mathrm{IC}\,\top} \mathbf{c}^{\mathrm{ret}}$
produces activations whose range is bounded; the softmax gain
$k > 0$ sharpens the distribution. Without $k$, the softmax
$P(j) \propto \exp(a_{j})$ over bounded activations is close to
uniform and the model fails to produce recency or contiguity
signatures. $k$ is fit alongside the other parameters; it is
initialized at the Cornell and Zhang (2025) Table 1 value
$k = 6.50$.

\paragraph{Primacy gradient.}
Standard CMR inherits a Polyn-et-al.-style primacy scaling
$p(t) = 1 + \phi \exp(-\psi t)$ applied to the
$M^{\mathrm{IC}}$ update magnitude at encoding step $t$. This is a
v6-native CMR mechanism, not an MS-TCM extension. Under the
\texttt{--standard-tcm} reduction the hierarchical layer is disabled
and the implementation falls back to a stronger primacy gradient
($\phi = 30.0$, $\psi = 0.8$) to keep the shape assertion satisfied;
under MS-TCM the per-storyline associative matrices contribute
additional primacy through the storyline-level beginning-of-list
synchronization at the start of each route-$\beta$ visit, so a
modest gradient ($\phi = 1.5$, $\psi = 0.5$) is sufficient.

\paragraph{No-repeats candidate set.}
Participants almost never repeat an item within one recall sequence
(intra-list repeat rates below $1\%$). A softmax readout without an
exclusion constraint, however, would predict the most recent item
with near-certainty whenever $k$ is moderate. We therefore mask the
softmax output to zero for any serial position already recalled on
the current list before re-normalising --- the standard CMR
convention. The likelihood routine skips observed repeats (rather
than returning $-\infty$) so the handful of real repeats in the
data do not poison the dataset log-likelihood.

\section*{S4. Numerical-accuracy anchors and model comparison}

The numerical anchor for MS-TCM's implementation has two parts.
First, the \citet{CornellZhang2025} hierarchical free-recall
implementation in \texttt{code/ms\_tcm/\_likelihood\_core.py} is
verified to $10^{-10}$ tolerance against hand-derived oracle tests at
\texttt{code/tests/test\_hcmr\_oracle.py}, so the comparator against
which MS-TCM's $\Delta\mathrm{LL}$ is reported is itself a faithful
implementation of its published inheritance. Second, when MS-TCM is
configured with $\lambda = 0$, $\tauinit = 0$, and a single storyline
($K = 1$), it reduces exactly to that C\&Z baseline at the same
machine-epsilon tolerance --- a regression covered by
\texttt{code/tests/test\_hcmr\_standard\_tcm.py}. The Layer 1 shape
assertions of \texttt{code/tests/test\_behavioral\_regression.py} pin
the canonical free-recall signatures (SPC, pFR, lag-CRP shapes) for
both the C\&Z baseline and MS-TCM. Together these anchors establish
that any reported $\Delta\mathrm{LL}$ between MS-TCM and the C\&Z
baseline is attributable to the new mechanisms (per-storyline
contexts, the global cross-storyline context, per-storyline
associative matrices, $\lambda$, the $\tauinit$-mixture between
recency-driven global recall and strict hierarchical recall) rather
than to implementation drift.

\newpage
\renewcommand{\refname}{Supplementary references}
\bibliographystyle{plainnat}
\bibliography{CDL-bibliography/cdl,local}

\end{document}
